\section{Tesis 3}

\subsection{Título}
\textit{Automatic Spanish Text Complexity Detection in Financial Documents with Low Data Availability} (Detección Automática de Complejidad de Texto en Español en Documentos Financieros con Baja Disponibilidad de Datos) \cite{romero2024}.

\subsection{Autor}
Mario Alberto Romero Sandoval, Instituto Tecnológico de Costa Rica.

\subsection{Resumen}
``El acceso a la información es un derecho humano fundamental en la sociedad moderna. Sin embargo, no todos tenemos el mismo acceso a la información, y una de las razones es que no comprendemos todo de la misma manera. El nivel educativo, la edad, las discapacidades y el contexto cultural pueden influir en la forma en que un texto es leído y comprendido por el público.

Ser capaz de discriminar entre segmentos de texto complejos y simples tiene numerosas aplicaciones: desde mejorar la eficiencia de los sistemas de simplificación, hasta aplicaciones educativas que ayudan a determinar si un texto es apropiado para el nivel de un estudiante determinado, y también a supervisar si las instituciones se están comunicando correctamente con el público.

En este trabajo se explorarán diferentes métodos y técnicas para la clasificación de textos en función de su complejidad, concretamente texto en español, así como métodos para resolver la escasez de datos en la tarea de discriminación de complejidad de texto en español. En particular, se centrará en el aprovechamiento de modelos de lenguaje existentes y el aprendizaje por transferencia para lograr y medir el impacto de datos aumentados mediante la generación de datos sintéticos en el problema de la discriminación de complejidad de texto.'' \cite{romero2024}

\subsection{Problema que busca resolver}
El campo de la simplificación automática de texto en español carece de conjuntos de datos etiquetados con pares de segmentos complejos y simples, los cuales son indispensables para entrenar clasificadores modernos de aprendizaje profundo. La mayoría de los enfoques existentes se centran en inglés u otros idiomas, y no existen modelos especializados para la clasificación de complejidad textual en español. Esta escasez de datos constituye un cuello de botella fundamental para el desarrollo de herramientas que permitan a poblaciones con comprensión lectora limitada (personas con impedimento visual, lectores no especializados) acceder a información financiera o técnica compleja. La pregunta central es si la generación automática de datos de simplificación mediante LLMs puede compensar la falta de datos etiquetados por humanos.

\subsection{Objetivos}

Objetivo General: Desarrollar y evaluar métodos de detección automática de complejidad de texto en español usando aprendizaje profundo y técnicas de aumentación de datos \cite{romero2024}.

Objetivos Específicos:
\begin{enumerate}
    \item Construir y/o identificar fuentes de datos para entrenar un modelo de aprendizaje profundo de detección de complejidad textual en español.
    \item Medir el impacto del tamaño del conjunto de datos y el uso de segmentos generados automáticamente en el problema de detección de complejidad textual en español.
    \item Proponer una metodología de aumentación de datos para un clasificador de segmentos de texto usando LLMs.
    \item Evaluar los efectos del ajuste fino (\textit{fine-tuning}) de modelos de clasificación con datos aumentados como forma de mejorar los resultados del aprendizaje por transferencia cuando se usan únicamente conjuntos de datos pequeños.
    \item Contrastar los resultados de usar modelos pre-entrenados como BETO/BERT como base para el aprendizaje por transferencia.
\end{enumerate}

\subsection{Hipótesis}

La tesis enuncia explícitamente dos hipótesis, lo cual es un aspecto metodológico destacable:

Hipótesis 1: \textit{Las capacidades de generación de texto de GPT-3 pueden utilizarse como fuente de texto aumentado para abordar la baja disponibilidad de datos de pares etiquetados de segmentos complejos y simples, proporcionando resultados similares al usarlo para ajustar modelos.}

Esta hipótesis fue rechazada: se encontró una diferencia estadística clara entre todas las fuentes de datos (valor-$p > 0.99$), confirmando que los datos manuales producen modelos significativamente superiores a los entrenados con datos generados por LLMs.

Hipótesis 2: \textit{El aprendizaje por transferencia en modelos pre-entrenados, usando datos aumentados, proporcionará mejoras en el rendimiento del modelo (F1-score) en comparación con modelos entrenados con datos pequeños pero generados manualmente.}

Esta hipótesis fue confirmada parcialmente: la aumentación de datos beneficia al modelo únicamente cuando el conjunto semilla es muy pequeño (aproximadamente 97 segmentos). Para conjuntos semilla de tamaño mayor, agregar datos aumentados no tiene efecto estadísticamente significativo o incluso degrada el rendimiento \cite{romero2024}.

\subsection{Resumen de la metodología empleada}

La metodología empleada por \cite{romero2024} se describe a continuación:

\begin{itemize}
    \item \textbf{Datos de partida:} Se usaron 5.314 pares de fragmentos de texto, cada par formado por una versión compleja y una simplificada del mismo contenido, extraídos de documentos de educación financiera y etiquetados manualmente por seis personas especializadas. Estos datos habían sido creados originalmente para apoyar a personas con impedimento visual.

    \item \textbf{Generación de datos adicionales artificiales:} Para estudiar si era posible compensar la escasez de datos reales, se generaron tres conjuntos alternativos de simplificaciones: uno usando ChatGPT sin instrucciones personalizadas, otro usando un modelo de lenguaje multilingüe llamado mT5, y un tercero usando un sistema automático de simplificación basado en reglas lingüísticas para idiomas románicos (Tuner).

    \item \textbf{Modelos evaluados:} Se compararon dos tipos de clasificadores: BETO, un modelo de inteligencia artificial preentrenado específicamente en texto en español que se ajustó para la tarea de detectar complejidad; y un clasificador más tradicional basado en estadísticas de frecuencia de palabras, usado como punto de referencia base.

    \item \textbf{Experimento 1:} Se compararon las cuatro fuentes de datos (manual, ChatGPT, mT5 y Tuner) con los dos modelos y veinte tamaños distintos de conjunto de entrenamiento, repitiendo cada combinación diez veces para obtener resultados estables. En total se realizaron 1.600 ejecuciones.

    \item \textbf{Experimento 2:} Se evaluó específicamente el efecto de agregar datos artificiales a distintos tamaños de datos reales, con 480 combinaciones distintas. Los resultados se validaron con una prueba estadística formal para confirmar que las diferencias observadas no fueran producto del azar.
\end{itemize}

\subsection{Principales resultados}

Los principales hallazgos reportados por \cite{romero2024} son:

\begin{itemize}
    \item \textbf{Los datos anotados por humanos producen mejores modelos:} Usando BETO, un modelo de inteligencia artificial especializado en español, los textos etiquetados manualmente obtuvieron una puntuación de precisión de 0.83 sobre 1. Los textos generados por ChatGPT alcanzaron 0.74, los de otro modelo de lenguaje (mT5) un 0.65, y los de un sistema automático de simplificación basado en reglas (Tuner) un 0.55. Las pruebas estadísticas confirmaron que la diferencia entre todas las fuentes es real y no producto del azar, rechazando así la primera hipótesis de la tesis.

    \item \textbf{El modelo deja de mejorar a partir de cierta cantidad de textos:} Todos los modelos llegan a un punto a partir del cual agregar más textos de entrenamiento ya no mejora su desempeño, alrededor de los 2.000 textos. El modelo entrenado con datos manuales continuó mejorando hasta los 4.000 textos aproximadamente, lo que indica que los datos humanos contienen información más rica.

    \item \textbf{Los datos artificiales solo ayudan cuando se tienen muy pocos datos reales:} Agregar textos generados por ChatGPT mejoró el rendimiento del modelo solo cuando el conjunto de datos reales era muy pequeño (alrededor de 97 textos). Con ese punto de partida y casi 500 textos artificiales adicionales, la mejora fue estadísticamente comprobable. Sin embargo, cuando ya se contaba con más de 492 textos reales, agregar datos artificiales empeoró el rendimiento.

    \item \textbf{La calidad de los datos importa más que la cantidad:} A mejor fuente de datos, mayor el rendimiento máximo alcanzable, independientemente del volumen generado. Producir más datos de baja calidad no compensa la falta de datos bien anotados por personas.
\end{itemize}

\subsection{Crítica de la tesis}
\begin{itemize}
    \item \textbf{Aspectos positivos:} La estructura de secciones es sólida y sigue una progresión metodológica bien articulada, desde la revisión de literatura hasta la validación experimental con pruebas estadísticas formales. A diferencia de otras tesis del área, las hipótesis se enuncian de forma explícita y verificable, y los objetivos específicos se corresponden directamente con las conclusiones: la Hipótesis 1 fue rechazada estadísticamente y la Hipótesis 2 fue confirmada parcialmente, con resultados cuantitativos claros para cada caso. El título es transparente sobre el alcance del trabajo al declarar desde el inicio las restricciones de dominio (documentos financieros) y de datos (baja disponibilidad), lo que es coherente con el trabajo relacionado que motivó la investigación.
    \item \textbf{Aspectos de posible mejora:} En primer lugar, los textos usados para evaluar el modelo son únicamente documentos financieros, lo que restringe la generalización de los resultados; el trabajo relacionado abarca dominios más amplios, por lo que el alcance del trabajo no queda totalmente alineado con la literatura que lo precede. En segundo lugar, GPT-3 se empleó sin instrucciones personalizadas ni las 21 reglas de simplificación usadas por los anotadores humanos, lo que hace que el rechazo de la Hipótesis 1 no sea concluyente sobre el potencial real de los LLMs con \textit{prompting} especializado. Por último, la tesis está escrita íntegramente en inglés para una institución costarricense cuyo dominio de aplicación es el español, sin que esta decisión se justifique en ninguna sección del documento.
\end{itemize}
